\section{Pseudo Code of Block Partitioning}
Algorithm~\ref{alg:extract_and_group} illustrates the block partitioning process and outputs a grouping map $G$, where each key corresponds to an ancestor node ID in the XML tree, and the associated value is a list of UI element indices that can be grouped together. Using these indices, we partition all the UI elements into $k = |G|$ UI blocks, denoted as $B = [b_1, b_2, ..., b_k]$.

\begin{algorithm}[H]
\caption{Layout-aware Block Partitioning}
\label{alg:extract_and_group}

\KwIn{Root element $e$ of an XML tree}
\KwOut{Grouping map $G$ from ancestor node ID to a list of important nodes' IDs}

\SetKwFunction{FExtract}{extract\_ancestor\_paths}
\SetKwFunction{FGroup}{group}
\SetKwProg{Fn}{Function}{:}{}
\BlankLine
\Fn{\FExtract{$e$, $A$, $M$}}{
    Append $e.index$ to $A$\;

    \If{$e$ is interactable and semantically meaningful}{
        $M[e.index] \gets A[:-1]$\;
    }

    \ForEach{child $c$ \textbf{in} $e.children$}{
        \FExtract{$c$, $A$, $M$}\;
    }

    Remove the last element from $A$\;
}
\BlankLine
\Fn{\FGroup{$M$}}{
    $L \gets$ maximum length of any value in $M$\;

    \For{$i \gets 0$ \KwTo $L - 1$}{
        $G \gets$ \{\}\;

        \ForEach{$(node, path)$ \textbf{in} $M$}{
            \eIf{$i < |path|$}{
                $g \gets path[i]$\;
            }{
                $g \gets path[-1]$\;
            }
            Append $node$ to $G[g]$\;
        }

        \If{$|G| \ge 3$}{
            \Return{$G$}
        }
    }

    \Return{$G$}
}
\BlankLine
\textbf{Main:}

$A \gets [\ ]$ \tcp*{List to store current ancestor path}

$M \gets \{\}$ \tcp*{Map from important node to its ancestor path}

\FExtract{$e$, $A$, $M$}\;

$G \gets \FGroup(M)$\;

\Return{$G$}
\end{algorithm}
